% Źródła: Eves s. 289--290 (zad. 10), 314--315; Wikipedia: Girard's theorem (Spherical trigonometry), Lexell's theorem.
% https://en.wikipedia.org/wiki/Lexell's_theorem
% https://en.wikipedia.org/wiki/Spherical_trigonometry

\section{Geometria sferyczna} % lang-pl
\section{Geometria sferica} % lang-it
\label{sekcja_sferyczna}

Druga geometria nieeuklidesowa jest starsza od pierwszej, tylko nikt jeszcze nie zauważy, że to geometria nieeuklidesowa: będzie to zwykła geometria na kuli. % lang-pl
W tej sekcji zobaczymy, że spełnia ona hipotezę kąta rozwartego, a potem poznamy jedno z jej twierdzeń, które przez sto lat zajmie największych geometrów. % lang-pl
La seconda geometria non euclidea è più antica della prima, solo che nessuno noterà ancora che è non euclidea: sarà la comune geometria sulla sfera. % lang-it
In questa sezione vedremo che soddisfa l'ipotesi dell'angolo ottuso e poi conosceremo un suo teorema che per cent'anni occuperà i più grandi geometri. % lang-it

Na sferze „prostymi'' są koła wielkie, dwie z nich zawsze się przecinają, a suma kątów trójkąta jest większa od $\pi$; Lambert zauważy tę analogię z hipotezą kąta rozwartego \cite[s. 286]{eves1_1972} (zobacz sekcję \ref{sekcja_lambert}). % lang-pl
Pole trójkąta sferycznego jest proporcjonalne do jego nadmiaru: dla sfery o promieniu $r$ i nadmiarze $E$ w stopniach jest ono równe $\pi r^2 E/180^\circ$ (Eves \cite[s. 289--290, zad. 10]{eves1_1972}). % lang-pl
Sulla sfera le «rette» sono i cerchi massimi, due di essi si intersecano sempre e la somma degli angoli di un triangolo è maggiore di $\pi$; Lambert noterà questa analogia con l'ipotesi dell'angolo ottuso \cite[p. 286]{eves1_1972} (si veda la sezione \ref{sekcja_lambert}). % lang-it
L'area di un triangolo sferico è proporzionale al suo eccesso: per una sfera di raggio $r$ ed eccesso $E$ in gradi vale $\pi r^2 E/180^\circ$ (Eves \cite[p. 289--290, es. 10]{eves1_1972}). % lang-it

\begin{theorem}
[Girarda] % lang-pl
[di Girard] % lang-it
\index{twierdzenie!Girarda} % lang-pl
\index{teorema!di Girard} % lang-it
	Pole trójkąta sferycznego o kątach $\alpha$, $\beta$, $\gamma$ na sferze o promieniu $R$ jest równe $R^2(\alpha + \beta + \gamma - \pi)$. % lang-pl
	L'area di un triangolo sferico con angoli $\alpha$, $\beta$, $\gamma$ su una sfera di raggio $R$ è uguale a $R^2(\alpha + \beta + \gamma - \pi)$. % lang-it
\end{theorem}

Albert Girard opublikuje je w 1629 roku; wcześniejszy, niepublikowany dowód znajdzie się w rękopisach Thomasa Harriota z 1603 roku. % lang-pl
Odpowiednikiem w geometrii hiperbolicznej będzie pole równe defektowi (zobacz sekcję \ref{sekcja_suma_katow}). % lang-pl
Albert Girard lo pubblicherà nel 1629; una dimostrazione precedente, inedita, si troverà nei manoscritti di Thomas Harriot del 1603. % lang-it
L'analogo nella geometria iperbolica sarà un'area uguale al difetto (si veda la sezione \ref{sekcja_suma_katow}). % lang-it
\index[persons]{Girard, Albert}%
\index[persons]{Harriot, Thomas}%

Eves \cite[s. 314--315]{eves1_1972} mówi krótko, że hipotezie kąta rozwartego odpowiada druga geometria nieeuklidesowa, że jest ona w pewnym sensie bardziej skomplikowana od geometrii Łobaczewskiego i że nie rozwija jej w swojej książce. % lang-pl
Eves \cite[p. 314--315]{eves1_1972} dice brevemente che all'ipotesi dell'angolo ottuso corrisponde una seconda geometria non euclidea, che per certi aspetti è più complicata di quella di Lobachevskij e che non la sviluppa nel suo libro. % lang-it
\todofoot{Geometria eliptyczna (sfera z utożsamionymi punktami antypodalnymi; nazwę wprowadzi Klein w 1871) oraz trygonometria sferyczna -- potrzebne źródło poza Evesem.} % lang-pl
\todofoot{Geometria ellittica (sfera con punti antipodali identificati; il nome lo introdurrà Klein nel 1871) e trigonometria sferica -- serve una fonte oltre Eves.} % lang-it

% Sferyczna osobno?
\begin{proposition}
[twierdzenie Lexella] % lang-pl
[teorema di Lexell] % lang-it
	Wierzchołki trójkątów sferycznych o ustalonej podstawie $AB$ i polu powierzchni leżą na okręgu małym (powstałym przez przecięcie sfery płaszczyzną) przechodzącym przez punkty antypodalne do wierzchołków $A$, $B$ z podstawy. % lang-pl
	I vertici dei triangoli sferici con base fissata $AB$ e area fissata giacciono su un cerchio minore (ottenuto intersecando la sfera con un piano) passante per i punti antipodali dei vertici $A$ e $B$ della base. % lang-it
\index{twierdzenie!Lexella} % lang-pl
\index{teorema!di Lexell} % lang-it
\end{proposition}

Anders Johan Lexell napisze około 1777 roku pracę na temat tego twierdzenia, razem z dwoma dowodami: trygonometrycznym i geometrycznym (wydrukowana będzie w 1784 roku). % lang-pl
Inni podadzą inne dowody: % lang-pl
Leonhard Euler w~1778 (kolega Lexella; jego dowód zostanie wydrukowany dopiero w 1797 roku), % lang-pl
Adrien-Marie Legendre w~1800, % lang-pl
Jakob Steiner w~1827, % lang-pl
Carl Friedrich Gauss w~1841, % lang-pl
Paul Serret w~1855, % lang-pl
Joseph-Émile Barbier w~1864. % lang-pl
Anders Johan Lexell scriverà intorno al 1777 un lavoro dedicato a questo teorema, contenente due dimostrazioni: una trigonometrica e una geometrica (sarà stampato nel 1784). % lang-it
Altri ne forniranno ulteriori dimostrazioni: % lang-it
Leonhard Euler nel 1778 (collega di Lexell; la sua dimostrazione sarà stampata solo nel 1797), % lang-it
Adrien-Marie Legendre nel 1800, % lang-it
Jakob Steiner nel 1827, % lang-it
Carl Friedrich Gauss nel 1841, % lang-it
Paul Serret nel 1855, % lang-it
Joseph-Émile Barbier nel 1864. % lang-it
\index[persons]{Lexell, Anders Johan}%
\index[persons]{Serret, Paul}%
\index[persons]{Barbier, Joseph-Émile}%
\index[persons]{Steiner, Jakob}%

Twierdzenie Lexella jest sferycznym odpowiednikiem twierdzeń I.37 i I.39 z \emph{Elementów}: trójkąty płaskie o wspólnej podstawie i równych polach mają wierzchołki na prostej równoległej do podstawy. % lang-pl
Il teorema di Lexell è l'analogo sferico delle proposizioni I.37 e I.39 degli \emph{Elementi}: i triangoli piani con base comune e aree uguali hanno i vertici su una retta parallela alla base. % lang-it

% Uwaga: dane o Lexellu pochodzą z Wikipedii (Lexell's theorem); Eves twierdzenia nie omawia.
% https://en.wikipedia.org/wiki/Spherical_trigonometry => Girard
